\scanpage{169}

\margin{Diviseur de fonctions}

Soit $X$ une variété irréductible \uline{normale} déf/$k$ $K=k(X)$ le corps des fonctions rationnelles/$k$ sur $X$

Soit $V$ une \deleted{\unsure{hypersurface irréductible}}\inserted{diviseur premier déf.} sur $k$ sur $X$

Alors l'ensemble des $f\in K$ \uline{définies} en $V$ forment un anneau de valuation discrète, $\deleted{\unsure{\underline O/V}}\underline O_{(V)}$.

En effet, c'est une \deleted{\illegible{1}} localité normale, dont la dimension de Krull est $1$.

Soit $\unsure{v|\nu}_V$ la valuation \inserted{de $K$} attachée à cet anneau de valuation du $K$. \unsure{Appelont|Appelons} \uline{diviseur} \inserted{déf/$k$} de $X$ toute comb. lin. formelle des diviseurs premiers \unsure{sur|rat.} déf/$k$. Posons, pour toute $f\in K$

\[
\mathrm{div}\,f=\sum_{\substack{V\ \mathrm{div.}\\ \mathrm{pr}/k}}v_V(f).\,V
\]

Cette somme a un sens, car pour $f$ donnée, il n'y a qu'un nb fini de termes non nuls ; \deleted{car \unsure{tout}} pour le voir, on est ramené au cas \deleted{\unsure{ouvert}}, affine, \unsure{alors|tous} les diviseurs pr. déf/$k$ correspondent aux idéaux premiers non nuls minimaux de l'anneau normal noethérien $A=k[X]$, et la chose est bien connue.

\uline{Théorème} \deleted{\unsure{pour} \illegible{1} \unsure{diviseur principal|diviseur équivalent}} \inserted{$\mathrm{div}\, f\geq 0$, $\Longleftrightarrow f$ est partout définie}

\unsure{Ceci|Cor.|C'est} résulte de la théorie des anneaux normaux : $A=\bigcap A_{\mathfrak p}$ ($\mathfrak p$ premier minimal) si $A$ est normal. En effet, on est aussitôt ramené au cas affine, et $\mathrm{div}\,f\geq0\Longleftrightarrow f\in\bigcap A_{\mathfrak p}$.

% archive sheet number: 168
